\section{周期化与边界处理}
\label{sec:wavelet-boundary}

Haar函数在$j\geq0$时支撑总在某个二进区间内，从不跨过$[0,1]$的端点，所以区间上的展开没有遇到边界问题。更光滑的紧支撑小波支撑更长，靠近端点的那些$\psi_{j,k}$有一部分伸到区间之外。要计算它们与$f$的内积，就得知道$f$在区间外取什么值，而区间外本来没有数据。怎样补上这些值，或者干脆改造边界附近的基函数，要由具体问题决定。本节先为周期问题构造周期基，再说明周期延拓对接缝处光滑性的影响，最后比较反射延拓与边界小波。

\subsection{周期化}
\label{sec:wavelet-periodization}

若要研究的对象本来就是周期$1$的函数，例如角度变量的函数，区间外的值由周期性给定，最自然的做法是把基函数也做成周期的。记$L^2(\mathbb T)$为周期$1$、在一个周期$[0,1)$上平方可积的函数（按几乎处处相等等同）组成的空间，内积取$[0,1)$上的积分。

\paragraph{周期化的定义}
设$\{V_j\},\varphi$是实值的正交多分辨率分析，$\varphi$与定理~\ref{thm:wavelet-mother-from-mra}的$\psi$都支撑在$[0,L]$中。对紧支撑函数$u$，令
\[
(\operatorname{Per}u)(x)=\sum_{n\in\mathbb Z}u(x+n).
\]
对每个$x$，只有有限个$n$使$u(x+n)\neq0$，所以和是有限和，$\operatorname{Per}u$是周期$1$的函数。对整数$j\geq0$和$0\leq k<2^j$，令\term{周期化}（periodization）后的基函数为
\[
\varphi^{\mathrm{per}}_{j,k}=\operatorname{Per}\varphi_{j,k},\qquad\psi^{\mathrm{per}}_{j,k}=\operatorname{Per}\psi_{j,k}.
\]
由$u_{j,k+2^jn}(x)=u_{j,k}(x-n)$，位置相差$2^j$的倍数的函数周期化后相同，所以每层只需取$0\leq k<2^j$。这里要求$j\geq0$，才能保证一个周期内的平移步长$2^{-j}$整除周期$1$。

\paragraph{展开公式}
计算内积要用一个简单的展开恒等式。设$u$有紧支撑且平方可积，$v$是周期$1$的平方可积函数，则
\begin{equation}
\int_0^1\overline{(\operatorname{Per}u)(x)}\,v(x)\,dx=\sum_{n\in\mathbb Z}\int_n^{n+1}\overline{u(y)}\,v(y-n)\,dy=\int_{\mathbb R}\overline{u(y)}\,v(y)\,dy .
\label{eq:wavelet-unfolding}
\end{equation}
第一个等号是有限和与积分交换并作代换$y=x+n$，第二个等号用到$v(y-n)=v(y)$。也就是说，周期函数与周期化函数在一个周期上的内积，等于原函数在全实轴上与周期函数的内积。

\begin{theorem}[周期小波基]
\label{thm:wavelet-periodic-basis}
在上述设定下固定整数$j_0\geq0$，则
\[
\{\varphi^{\mathrm{per}}_{j_0,k}:0\leq k<2^{j_0}\}\cup\{\psi^{\mathrm{per}}_{j,k}:j\geq j_0,\ 0\leq k<2^j\}
\]
是$L^2(\mathbb T)$的标准正交基。于是每个$f\in L^2(\mathbb T)$都按这组基在$L^2(\mathbb T)$中展开，系数的平方和等于$\int_0^1\lvert f\rvert^2$。
\end{theorem}

\begin{proof}
\emph{正交性与归一化。}取$j,j'\geq j_0$及相应的$k,k'$。由式~\eqref{eq:wavelet-unfolding}，
\[
\langle\psi^{\mathrm{per}}_{j,k},\psi^{\mathrm{per}}_{j',k'}\rangle_{\mathbb T}=\int_{\mathbb R}\psi_{j,k}\,\psi^{\mathrm{per}}_{j',k'}\,dy=\sum_{n\in\mathbb Z}\langle\psi_{j,k},\psi_{j',k'-2^{j'}n}\rangle .
\]
由定理~\ref{thm:wavelet-real-line-expansion}，右端每一项在$j=j'$且$k=k'-2^{j'}n$时为$1$，否则为$0$。$0\leq k,k'<2^j$时，只有$n=0$且$k=k'$才可能。所以这些函数标准正交。$\varphi^{\mathrm{per}}_{j_0,k}$之间、以及它们与$\psi^{\mathrm{per}}_{j,k}$之间的内积用同样的计算，并用$V_{j_0}\perp W_j$（$j\geq j_0$）。

\emph{完备性。}记$V^{\mathrm{per}}_J$为$\{\varphi^{\mathrm{per}}_{J,k}\}_{0\leq k<2^J}$的线性张成。给定$g\in L^2(\mathbb T)$，令$f=g\mathbf{1}_{[0,1)}\in L^2(\mathbb R)$，则$\operatorname{Per}f=g$。$\varphi_{J,k}$支撑在$[k2^{-J},(k+L)2^{-J}]$中，只有$-L<k<2^J$时才与$[0,1)$相交，所以$P_Jf$是有限和，且对一切$J\geq0$都支撑在$[-L,1+L]$中。对支撑在这个区间中的$u$，每个$x$处至多有$2L+2$个非零项$u(x+n)$，由Cauchy--Schwarz不等式，
\[
\int_0^1\lvert\operatorname{Per}u\rvert^2\leq(2L+2)\int_{\mathbb R}\lvert u\rvert^2 .
\]
于是$\lVert g-\operatorname{Per}P_Jf\rVert_{\mathbb T}^2\leq(2L+2)\lVert f-P_Jf\rVert_2^2\to0$。而$\operatorname{Per}P_Jf$是有限个$\operatorname{Per}\varphi_{J,k}$的组合，属于$V^{\mathrm{per}}_J$。所以$\bigcup_JV^{\mathrm{per}}_J$在$L^2(\mathbb T)$中稠密。

最后核对每个$V^{\mathrm{per}}_J$都被所给函数族张成。紧支撑时二尺度系数只有有限个非零，$\varphi_{j_0,k}$和$\psi_{j,k}$（$j<J$）都是有限个$\varphi_{J,n}$的组合，周期化后属于$V^{\mathrm{per}}_J$。所给族中$j<J$的部分有$2^{j_0}+2^{j_0}+\cdots+2^{J-1}=2^J$个标准正交函数，而$V^{\mathrm{per}}_J$由$2^J$个函数张成，维数不超过$2^J$，所以它们恰好构成$V^{\mathrm{per}}_J$的基。由稠密性，整个族是完备的，展开与Parseval等式由第~\ref{chap:inner-product-spaces}章得到。
\end{proof}

对Haar，$j\geq0$时$\psi_{j,k}$的支撑在$[0,1)$内，周期化不改变它们在一个周期上的取值，定理~\ref{thm:wavelet-periodic-basis}回到定理~\ref{thm:wavelet-haar-expansion}。

\paragraph{消失矩只部分保留}
零阶矩保留：由式~\eqref{eq:wavelet-unfolding}取$v\equiv1$，$\int_0^1\psi^{\mathrm{per}}_{j,k}=\int_{\mathbb R}\psi_{j,k}=0$。高阶矩却不能照搬。设$\psi_{j,k}$的支撑在$(0,2)$中且跨过$x=1$，在$[0,1)$上计算，
\begin{equation*}
\begin{aligned}
\int_0^1x\,\psi^{\mathrm{per}}_{j,k}(x)\,dx&=\int_0^1y\,\psi_{j,k}(y)\,dy+\int_1^2(y-1)\,\psi_{j,k}(y)\,dy\\
&=\int_{\mathbb R}y\,\psi_{j,k}(y)\,dy-\int_1^{\infty}\psi_{j,k}(y)\,dy .
\end{aligned}
\end{equation*}
即使$\psi$有两个消失矩使第一项为零，第二项一般也不为零。原因在于：把$x$看成周期函数时，它在整数点处从$1$跳回$0$，在跨过接缝的支撑上已不是多项式。全实轴上“与多项式内积为零”的性质，只对支撑不跨接缝的那些周期小波原样成立。

\subsection{周期延拓的边界效应}
\label{sec:wavelet-periodic-extension}

上一小节的基是为周期函数准备的。实际问题常常是给定$[0,1)$上的函数$f$，把它按周期$1$重复延拓成周期函数$\tilde f(x)=f(x-\lfloor x\rfloor)$，再用周期基展开。反过来，周期函数限制在$[0,1)$上也得到区间上的函数。两个方向的映射互逆，并且$\int_0^1\lvert\tilde f\rvert^2=\int_0^1\lvert f\rvert^2$，所以\term{周期延拓}（periodic extension）是$L^2([0,1))$到$L^2(\mathbb T)$的线性等距同构。这个对应只涉及一个周期上的积分，不要求$f(0)$与$f(1)$有任何关系；事实上在$L^2$中单点取值根本没有意义。

\paragraph{区间上光滑不等于周期光滑}
$L^2$等距不保证光滑性被保留。取$f(x)=x$，它在$[0,1]$上无穷次可微，周期延拓却是锯齿波，在每个整数点处从接近$1$跳回$0$。这个跳跃在系数中看得见。用周期$1$的Fourier系数$\hat f(n)=\int_0^1e^{-2\pi\mathrm inx}f(x)\,dx$计算，分部积分得
\[
\hat f(n)=\Bigl[\frac{x\,e^{-2\pi\mathrm inx}}{-2\pi\mathrm in}\Bigr]_0^1+\frac{1}{2\pi\mathrm in}\int_0^1e^{-2\pi\mathrm inx}\,dx=\frac{\mathrm i}{2\pi n}\qquad(n\neq0),
\]
只按$1/\lvert n\rvert$衰减，这是带跳跃的周期函数的典型速度。在小波系数中，跨过接缝的那些$\psi^{\mathrm{per}}_{j,k}$看到的是跳跃而不是光滑的直线，定理~\ref{thm:wavelet-moment-decay}的估计对它们不再适用，只剩下$\lvert\langle\psi^{\mathrm{per}}_{j,k},\tilde f\rangle\rvert\leq\sup\lvert f\rvert\cdot\lVert\psi_{j,k}\rVert_1=\sup\lvert f\rvert\cdot2^{-j/2}\lVert\psi\rVert_1$这样的粗略界。每层跨接缝的小波只有有限个（不超过$L+1$个），但它们的系数衰减慢，正是边界效应的来源。

\paragraph{端点相容条件}
对在$[0,1]$上连续的$f$，周期延拓连续当且仅当$f(0)=f(1)$；若$f$在$[0,1]$上$m$次连续可微，延拓$m$次连续可微当且仅当$f^{(i)}(0)=f^{(i)}(1)$对$i=0,1,\ldots,m$都成立。这些条件要逐阶分别核对。例如$f(x)=x(1-x)$满足$f(0)=f(1)=0$，延拓是连续的；但$f'(0)=1\neq-1=f'(1)$，延拓在整数点处有折角。它的Fourier系数可由前面的计算得到：$\int_0^1x^2e^{-2\pi\mathrm inx}\,dx=\mathrm i/(2\pi n)+1/(2\pi^2n^2)$，所以
\[
\widehat{x(1-x)}(n)=\frac{\mathrm i}{2\pi n}-\Bigl(\frac{\mathrm i}{2\pi n}+\frac{1}{2\pi^2n^2}\Bigr)=-\frac{1}{2\pi^2n^2}\qquad(n\neq0),
\]
按$1/n^2$衰减，比锯齿波快，又因为导数在接缝处跳跃而不能更快。周期延拓的接缝正则性就是这样一阶一阶决定系数行为的。

\subsection{区间边界的处理}
\label{sec:wavelet-interval-boundary}

若函数本来只定义在$[0,1]$上，端点处的行为有实际意义，例如边界上的取值或流量，那么人为引入接缝跳跃会让边界附近的系数失真。常见的处理有两种思路：一是换一种延拓方式，给出更合理的区间外取值；二是不延拓，直接在区间内重新构造靠近端点的基函数。

\paragraph{反射延拓}
把$f$先关于$x=0$作偶反射，令$F(x)=f(\lvert x\rvert)$（$-1\leq x\leq1$），再按周期$2$延拓。这称为\term{反射延拓}（reflection extension）。若$f$在$[0,1]$上连续，$F$在$x=0$和$x=\pm1$处左右取值自动相等，延拓是连续的，取值跳跃被消除了。导数却不一定连续：$F'(0^+)=f'(0)$而$F'(0^-)=-f'(0)$，只有$f'(0)=0$时才连续，在$x=1$处同理。

仍以$f(x)=x$为例，反射延拓是周期$2$的三角波$\lvert x\rvert$。按周期$2$的Fourier系数计算，
\[
\frac12\int_{-1}^1\lvert x\rvert e^{-\mathrm i\pi nx}\,dx=\int_0^1x\cos(\pi nx)\,dx=\frac{(-1)^n-1}{\pi^2n^2}\qquad(n\neq0),
\]
奇数$n$时为$-2/(\pi^2n^2)$。与周期延拓的$1/\lvert n\rvert$相比，衰减快了一阶，但接缝处的折角使它停在$1/n^2$。所以反射可以消除取值跳跃，却不自动保留高阶导数的连续性；用反射延拓后的函数再作周期小波展开时，接缝处的系数仍受折角影响。

\paragraph{边界小波}
另一种思路是保留那些支撑完全落在$[0,1]$内部的$\psi_{j,k}$，把跨过端点的那些替换为每端有限个专门构造的\term{边界小波}（boundary wavelets）。构造的目标是：限制在$[0,1]$上的、次数小于$r$的多项式仍属于各层尺度空间，从而边界附近的细节系数也能抵消低次多项式，定理~\ref{thm:wavelet-moment-decay}式的衰减在端点附近同样成立。

调整后的函数族不再由单个函数平移缩放得到，前面的定理都不能直接引用，必须重新核对几件事：边界函数与内部函数之间的正交关系，或者在放弃正交时与另一组对偶函数的对偶关系；整个族在$L^2([0,1])$中的完备性；以及边界函数实际具有的消失矩个数。这些构造的细节和对应的有限滤波实现，属于信号处理中的计算问题，可参考\BookChapterRef{chap:signal-analysis}{卷九第6章“信号的表示与比较”}。

三种处理各有适用场合：周期化适合本来就周期的对象；反射延拓简单，能消除取值跳跃，适合只需要连续性的场合；边界小波保留端点处的高阶正则性，代价是构造复杂。选哪一种，取决于区间外取值在问题中意味着什么。
